\subsubsection{设置对照方法并建立强化学习调度模型}\label{sec:q3-7}

四类方法共同遵守设备反馈、几何清除和完整退出约束，区别在于如何选择下一项行动。
前瞻基线先生成少量候选任务，在近似场景下展开有限尾程，依据预计总费用选择首动作；
它已经包含信息价值近似，因而是一个较强对照。
几何法使用最近邻与2-opt安排全局任务次序，以有限几何假设和误差样本估计主动探测后的定位费用，并允许覆盖点在保持全覆盖的条件下移动。
状态搜索则使用第\ref{sec:q3-3}节的显式任务状态与下界剪枝。
三者分别体现有限前瞻、几何局部优化和结构化组合规划。

强化学习采用PPO策略优化\cite{schulman2017ppo}。
网络对每个合法候选的特征与全局上下文编码，输出动作概率；
候选覆盖测点、频道、是否中断扫描、待处理源和已有清除证书等选择。
候选集合随观测更新，非法动作通过掩码排除。
几何区域、有限候选生成及有界退出仍由规则维护，网络在这些约束下学习调度，而非直接回归任意连续坐标。
首版最终模型使用96维隐藏层，共62882个参数。

令策略参数为 \(\phi\)，某步题目计费用时增量为 \(\Delta T_t\)，常规奖励为 \(r_t=-\Delta T_t/1000\)。
固定兜底费用与失败罚项也计入训练回报。
采用终止式、无时间折扣的目标 \(\gamma=1\)，使累计时间成本与第三问任务目标一致。
PPO通过概率比 \(w_t(\phi)=\pi_\phi(a_t\mid S_t)/\pi_{\rm old}(a_t\mid S_t)\) 构造裁剪目标

\begin{equation}\label{eq:q3-15}
\mathcal L_{\rm clip}(\phi)=\mathbb E_t\left[
\min\{w_t(\phi)\widehat A_t,\ 
\operatorname{clip}(w_t(\phi),1-\eta,1+\eta)\widehat A_t\}
\right].
\end{equation}

其中 \(\eta\) 是策略比率裁剪宽度；
价值网络提供回报基准，采用广义优势估计\cite{schulman2016gae}

\begin{equation}\label{eq:q3-16}
\widehat A_t=\sum_{\ell=0}^{H-t-1}\lambda^\ell\delta_{t+\ell},
\qquad \delta_t=r_t+V(S_{t+1})-V(S_t),
\end{equation}

其中 \(H\) 为该条终止轨迹的动作步数，\(\lambda\) 为GAE系数，首版最终模型取0.95；
终止状态的后续价值取零。
有限候选减轻了连续动作探索的难度，也限定了学习策略能够发现的动作范围。
因此，我们将强化学习视为检验``在相同任务结构下，学习调度能否改善人工规划''的独立方向，并通过等预算配置比较和跨场景评估判断效果。

\subsubsection{设置实验场景与评价指标}\label{sec:q3-8}

为检验几何规则、显式路线规划与学习调度在第三问中的作用，我们比较冻结的前瞻 rollout 基线、状态搜索、PPO 强化学习和几何探测四种方法。
以前瞻策略作为较强参照；
状态搜索与几何法均包含经覆盖条件认证的未来站点移动，PPO 则在固定七个覆盖点及合法定位候选中选择动作。
因此，主实验比较的是四套可执行方案的综合效果，单个模块的贡献另由消融实验分析。

首版实验在同一 Linux 主机、相同场景生成规则和逐场景噪声场上进行。
先用开发集诊断，以独立的 48 局筛选候选，冻结算法、参数和模型后，才运行 256 个随机场景及 28 个压力场景。
研究分布将源数设为 10---16，源在半径 1800 m 圆域内按面积均匀分布，接收半径在 1000---1500 m 间均匀抽样；
压力集包含七类结构、每类四局，单独报告。
以全清局数占总场景数的比例和失败清除次数检查可行性，以整局平均耗时、P95 及逐场景差值衡量效率。
平均节省的 95\% 区间由 10000 次配对 bootstrap 给出，正值表示候选较参照节省时间。
按场景报告配对区间，有助于避免仅依赖点估计进行比较\cite{agarwal2021}。
该区间描述冻结策略面对场景抽样的变动，并不涵盖重新训练网络的随机性。

为便于区分指标，令第 \(i\) 个场景的基线与候选耗时分别为 \(T_i^{(b)}\) 和 \(T_i^{(m)}\)，定义

\begin{equation}\label{eq:q3-17}
\overline T_m=\frac1n\sum_{i=1}^nT_i^{(m)},\qquad
\overline\Delta_m=\frac1n\sum_{i=1}^n(T_i^{(b)}-T_i^{(m)}),\qquad
G_m=\frac{\overline\Delta_m}{\overline T_b}.
\end{equation}

全清失败或超时必须计入失败记录，不能只对成功子集报告速度优势。
随机集与压力集分别统计全清比例，两种分布的时间统计也分别计算。

\Needspace{82mm}
